\documentclass[11pt]{article}

\usepackage[margin=1in]{geometry}
\usepackage{amsmath,amssymb,amsthm}
\usepackage{graphicx}
\usepackage{hyperref}
\usepackage{xurl}
\usepackage{microtype}
\usepackage{enumitem}
\usepackage{booktabs}
\usepackage{tabularx}
\usepackage{array}
\newcolumntype{Y}{>{\raggedright\arraybackslash}X}
\usepackage{listings}

\graphicspath{{../figures/}}

\newtheorem{theorem}{Theorem}
\newtheorem{lemma}{Lemma}
\newtheorem{proposition}{Proposition}
\newtheorem{corollary}{Corollary}
\theoremstyle{definition}
\newtheorem{definition}{Definition}
\theoremstyle{remark}
\newtheorem{remark}{Remark}

\newcommand{\TV}{\mathrm{TV}}
\newcommand{\eps}{\varepsilon}

\title{PSC-Q: Public Service Calendars and Slack-Harvesting Substitution Padding \\ for Congestion-Private Anonymous Overlays}
\author{}
\date{Draft v0.33 (2026-03-17) (archive v290)}



\begin{document}
\maketitle

\begin{abstract}
We study practical mechanisms for meeting Congestion-EQ-style budgets in anonymous overlays, with emphasis on Anonymous-DHT lookups.
A central engineering pitfall is that \emph{additive} padding can worsen congestion leakage by pushing the system toward high utilization.
We develop \emph{PSC-Q} (Public Service Calendar for Queues): a mechanism family that (i) publishes a service calendar, (ii) harvests slack via \emph{substitution} (idle-filling) rather than additive traffic, and (iii) uses declared dither to contract TV and to attenuate periodic witnesses.
We connect PSC-Q to a knob-synthesis view (including slack-aware expiry synthesis, SAES) and provide new evidence quantifying the discretization ``$\delta$-tax'' for periodic calendars.
The contribution is deliberately narrow: PSC-Q is the complete practical mechanism companion beneath Congestion-EQ, and it owns only the declared \emph{mechanism packet} (calendar family, substitution semantics, dither/expiry knobs, utilization range, observer assumptions, and minimal calibration evidence) that explains how a deployment intends to meet a Congestion-EQ budget.
Witness/accountant semantics remain with Congestion-EQ, the supporting-evidence packet remains with the PSC-Q addendum, and replayable verifier bundles, notarized evaluation logs, and public release lineage are imported from W-Congestion-EQ, Eval~1, and Release~A rather than re-owned here. That division makes PSC-Q referee-complete as the congestion-family mechanism paper without turning it into a second witness/accountant, evaluator, or release note.
\end{abstract}
\paragraph{Compilation note.}
This paper is brief-by-design. A companion addendum collects the operator-facing microstudies and checklists that were previously shipped as a split supplement.\cite{series:pscqaddendum}

\paragraph{Scope.}
We develop practical compiler mechanisms to meet Congestion-EQ budgets without shrinking slack: public service calendars, slack-harvesting \emph{substitution} padding, declared dither/expiry knobs, and the compact mechanism packet that records those choices.
This paper is intentionally the downstream mechanism companion rather than the congestion-family entrypoint, and it stops before the supporting-evidence addendum, congestion-specific verifier bundles, generic notarized evaluation logs, runtime conformance receipts, or public release lineage.\cite{series:pscqaddendum,series:wcongeq,series:evalnotary,series:opb,anondht_release_receipts}

\paragraph{Imports.}
We import the Congestion-EQ contract (Paper~I) as the target guarantee.\cite{series:congeq}
We also import the series' destination/stop-time compiler interface and TV certificate vocabulary.\cite{series:schedcompiler,series:knapsacktransport}
Verifier-bundle replay, evaluator notarization, and public release lineage are explicitly treated as imported companion layers rather than PSC-Q responsibilities.\cite{series:wcongeq,series:evalnotary,anondht_release_receipts}

\paragraph{Exports.}
This paper exports a declared \emph{PSC-Q mechanism packet}: calendar family and service model, substitution policy, dither/expiry knobs, utilization range, observer assumptions, and the minimal calibration evidence needed to justify those knobs.
SAES appears only as a knob-synthesis rule for choosing that packet; later verifier/notarization/release layers consume the packet but are not redefined here. In particular, PSC-Q is intended to be the first public congestion-family mechanism note a referee can read after Congestion-EQ, with the addendum and W-Congestion-EQ available only if the audit question later widens to operator microstudies or replayable checker bundles.

\paragraph{Boundary and earliest sufficient stop.}
This paper is complete once an auditor can identify the declared calendar family, substitution semantics, dither/expiry knobs, utilization range, observer assumptions, and the minimal calibration evidence that a compiler used to size them.
That is enough to read PSC-Q as the practical congestion-family companion beneath Congestion-EQ. If the next question is ``does this exact packet replay under a checker?'', that handoff belongs to W-Congestion-EQ; if the next question is ``what additional operator microstudies or checklists support these knobs?'', that handoff belongs to the PSC-Q addendum; if the next question is ``was the evaluation workflow notarized without grinding?'', that handoff belongs to Eval~1; if the next question is ``how does this packet sit inside a public release lineage?'', that handoff belongs to Release~A.\cite{series:pscqaddendum,series:wcongeq,series:evalnotary,anondht_release_receipts}

\paragraph{Artifact policy.}
This archive is source-only; supporting plots and scripts are available on request.\cite{series:artifactpolicy}

\paragraph{Later-terse-paper citation posture.}
When a later short paper only needs to answer \emph{which public service calendar / substitution / dither packet was declared for this congestion-private anonymous-DHT family?}, it should stop at the minimal public PSC-Q mechanism card below plus the tiny slot-budget drill.
It should widen only when the live issue becomes the congestion witness/accountant root (Congestion~1), optional operator microstudies (PSC-Q Addendum), replayable congestion-specific checking (W-Congestion-EQ), evaluator notarization (Eval~1), runtime plan conformance (Operational~B), or public release lineage (Release~A).\cite{series:congeq,series:pscqaddendum,series:wcongeq,series:evalnotary,series:opb,anondht_release_receipts}

\begin{table}[t]
\centering
\small
\begin{tabularx}{0.97\linewidth}{@{}p{0.30\linewidth}Y@{}}
\toprule
\textbf{Public row} & \textbf{Pinned value}\\
\midrule
Mechanism card id & \texttt{calendar\_manifest\_id = pscq.knobs.v0.1}; this card is the public mechanism packet, not the upstream congestion witness/accountant object \\
Calendar family & \texttt{renewal\_uniform} with public tick width $\delta = 2$ ms (representative feasible row from Table~\ref{tab:pscq_knobs}) \\
Substitution semantics & Cover work fills only service opportunities that would otherwise be idle; real work preempts cover so the packet does not claim additive padding on top of a busy queue \\
Dither / observer contract & Declared public jitter is sized against the observer-resolution test $(1/2-\alpha)\delta \lesssim \sigma_{\mathrm{obs}}$; the packet therefore commits to an observer-noise assumption rather than merely naming a calendar family \\
Utilization range & Compile for $\rho \in [0.7,0.8]$ and treat $\rho_{\max}=0.8$ as part of the public privacy contract; recertify if runtime monitoring leaves that range \\
Minimal calibration evidence & Publish the knob-table digest and the supporting sweep row(s) needed to justify the chosen $(\delta,\alpha)$ / utilization packet; optional wider operator microstudies stay in the addendum \\
Owner routing & Congestion-EQ owns witness/accountant semantics; W-Congestion-EQ owns replayable checker bundles; Eval~1 / Operational~B / Release~A own notarization, runtime conformance, and lineage-tagged release receipts \\
\bottomrule
\end{tabularx}
\caption{A minimal public PSC-Q mechanism card. This is the smallest public answer later terse papers can quote directly when the live issue is the declared congestion-family mechanism packet rather than the witness/accountant root, optional operator evidence, checker replay, evaluator logs, runtime conformance, or release lineage.}
\label{tab:minimal-pscq-card}
\end{table}

This card is intentionally non-decorative: it pins one calendar family, one substitution rule, one utilization interval, one observer-side resolution contract, and one minimal evidence obligation, without silently widening into the congestion theorem root or later checker / release layers.

\paragraph{One tiny slot-budget drill.}
If the declared calendar tick is $\delta=2$ ms, the public service calendar offers
\[
\frac{1}{\delta} = \frac{1}{0.002\ \mathrm{s}} = 500
\]
service opportunities per second.
At the declared utilization ceiling $\rho_{\max}=0.8$, the packet therefore leaves an average slack budget of
\[
(1-\rho_{\max})\cdot 500 = 0.2\cdot 500 = 100
\]
otherwise-idle opportunities per second for substitution-style cover.
A later terse paper may therefore say ``the PSC-Q packet declared a 2 ms public service calendar with about 100 slack opportunities per second at the boundary utilization regime'' without reopening the operator microstudy tail or the release stack.


\section{Introduction}
PSC-Q is the downstream mechanisms companion to the Congestion-EQ contract \cite{series:congeq}.
Where Congestion-EQ defines \emph{what} must be bounded (a congestion-witness family plus an explicit accountant for repeats/resets), PSC-Q focuses on \emph{how a compiler can actually meet those budgets} in a deployed overlay or Anonymous-DHT lookup service without turning the theorem root itself into a compiler-design note.

\paragraph{Series map.}
We assume the series' lookup-compiler interface for destination/contact traces and stop times \cite{series:schedcompiler,series:knapsacktransport}.
Congestion-EQ \cite{series:congeq} introduces congestion as a first-class leakage tier.
This paper proposes public-service-calendar mechanisms and slack-harvesting \emph{substitution} padding to keep utilization in a safe regime.
It is intended to be the stable practical follow-on to Congestion-EQ: once the theorem/accountant root is accepted, PSC-Q supplies the mechanism packet a deployment can actually declare before any verifier bundle, evaluator notarization, or release-lineage packaging is considered. The assurance companion \cite{series:wcongeq} shows how to bind those knob choices into a replayable congestion-specific verifier bundle for later evaluation/release layers; PSC-Q itself stops at the declared mechanism packet rather than re-owning that verifier surface.

This paper is intentionally mechanisms-first.
For motivation, threat surface, and the broader congestion/latency side-channel literature, see Congestion-EQ \cite{series:congeq}.
Here we focus on the compiler-facing issue: \emph{padding that increases offered load can worsen congestion leakage} by shrinking slack.
PSC-Q therefore emphasizes \emph{substitution} (idle-filling) semantics and explicit slack targets; the PSC-Q addendum remains only optional supporting evidence and W-Congestion-EQ remains only the later replayable checker layer.\cite{series:pscqaddendum,series:wcongeq}

\paragraph{Earliest practical stop.}
A referee or operator can stop after this paper once the congestion-family question becomes \emph{``what calendar / substitution / dither packet should a deployment declare?''}
That is the point of PSC-Q in the stack.
Only narrower later questions---``show me the extra microstudies,'' ``show me the replayable checker bundle,'' ``show me the notarized evaluator transcript,'' or ``show me the public release lineage''---require the addendum, W-Congestion-EQ, Eval~1, or Release~A.\cite{series:pscqaddendum,series:wcongeq,series:evalnotary,anondht_release_receipts}

\paragraph{Anonymous-DHT burstiness (pointer).}
Anonymous-DHT lookups can be bursty because they combine bounded parallelism, retries, and provider discovery.
Rather than restate deployed-DHT measurement context, we point to Congestion-EQ \S1 for the Anonymous-DHT motivation and citations, and to the Boss Fight dial-sheet paper for knob-level compilation context.


\section{Padding taxonomy: additive vs substitution}
\begin{definition}[Additive vs substitution padding]
\emph{Additive} padding injects extra work on top of real work.
\emph{Substitution} padding (idle-filling) uses otherwise-idle capacity to serve cover traffic in place of idle time, without increasing mean load.
\end{definition}

\begin{remark}[Substitution must not steal slack]
``Substitution'' is slack-harvesting only if cover work is strictly limited to otherwise-idle service opportunities (e.g., calendar slots that would have been idle).
If cover traffic competes with real work in the same queue (or contends for shared downstream resources), it can bias service and increase delay, defeating the congestion-privacy goal.
Section~\ref{sec:microstudies-summary} includes a small service-bias microstudy and an implementation checklist for ``preemptible'' cover.
\end{remark}

The substitution/idle-filling idea mirrors classic constant-time countermeasures in software: when a secret-dependent branch would ``finish early,'' the implementation executes dummy work to fill the slack and equalize an adversary-visible timing trace \cite{agat2000,mishra_fallthrough2025}.

Substitution-style padding is deployed in Tor as hop-local connection- and circuit-level padding machines (negotiation in Proposal~254 and the Tor padding specification) \cite{torproposal254,torpaddingspec}.
By contrast, classic website-fingerprinting defenses such as BuFLO and CS-BuFLO primarily add work to approach a constant-rate trace \cite{dyer2012peekaboo,csbuflo}.
This taxonomy matters for congestion privacy: additive padding can erode slack and worsen Congestion-EQ witnesses even when the padding is secret-independent.

The archive contains direct evidence that additive padding can increase distinguishability by shrinking slack.

\begin{figure}[t]
\centering
\fbox{\begin{minipage}{0.92\linewidth}\centering\small Figure omitted in source-only release. Requested file: \texttt{fig\_additive\_vs\_substitution\_tv.pdf}.\end{minipage}}
\caption{Archive evidence: substitution padding can reduce TV while additive padding can increase it at the same nominal ``cover rate.''}
\label{fig:tax}
\end{figure}

\section{PSC-Q: Public service calendars}
PSC-Q fixes a public calendar of service opportunities (cells, quanta, or send slots) and then uses substitution semantics to fill idle opportunities with cover traffic.
The calendar itself is a public artifact and can be audited.

\subsection{Why public calendars help}
A public calendar separates two concerns:
(i) schedule regularity and capacity provisioning are public and auditable;
(ii) secret-dependent variability is forced into controlled degrees of freedom (e.g., which cover items occupy idle slots), enabling clean accounting.

\subsection{Idle filling and partial-idle semantics}
A PSC-Q calendar can be instantiated as a stream of public service opportunities (continuous-time rate $\mu$ or periodic ticks of width $\delta$).
At each opportunity, the relay serves one real cell if available; otherwise it emits a \emph{hop-local} dummy that terminates at the next hop.
If \emph{every} idle opportunity is filled, the link-level departure timestamps match the public calendar exactly, eliminating hop-local rate/count leakage.
Crucially, because dummies consume only idle service, the real-work queue dynamics are unchanged: this mechanism helps the \emph{trace tier} without automatically fixing the \emph{congestion tier} (Congestion-EQ).

Many deployments intentionally fill only a fraction of idle opportunities to save bandwidth.
Partial idle filling reduces but does not eliminate departure-count distinguishability; Section~\ref{sec:microstudies-summary} summarizes this residual leak and the resulting overhead tradeoffs.

\begin{lemma}[Slack-harvesting semantics do not change the real queue]
Consider a single-server queue fed by \emph{real} arrivals, with a public PSC-Q service calendar specifying service opportunities.
At each opportunity the server serves one real job if available; otherwise it serves a hop-local dummy that consumes only that opportunity.
Then (i) the real-queue length process and real waiting-time distribution are identical to the baseline system that simply idles when the real queue is empty, and (ii) the departure timestamps of \emph{served slots} match the public calendar.
\end{lemma}

\begin{proof}
Couple the PSC-Q system to the baseline on the same real-arrival sample path and the same calendar.
Both systems serve a real job whenever one is present at a service opportunity; PSC-Q differs only by emitting a dummy when the real queue is empty.
Thus the two systems have identical real departures and queue evolution, while PSC-Q replaces baseline idle slots with hop-local dummy service.
\end{proof}

\section{Multi-hop effects: bottleneck shift and localization}
Even if hop-local congestion is controlled, multi-hop overlays can exhibit bottleneck shift: a defense that calms one queue can move congestion elsewhere.

\begin{figure}[t]
\centering
\fbox{\begin{minipage}{0.92\linewidth}\centering\small Figure omitted in source-only release. Requested file: \texttt{fig\_tandem\_padding\_bottleneck\_tv.pdf}.\end{minipage}}
\caption{Archive evidence: tandem/multi-hop systems can shift leakage to a different hop under naive padding. Mechanisms should track a witness family, not a single queue.}
\label{fig:tandem}
\end{figure}

\section{Knob synthesis: slack-aware expiry synthesis (SAES)}
PSC-Q exposes knobs: calendar period $\delta$, flush strength, cover pool size, dither distribution, and expiry rules.
The archive develops SAES as a helper that chooses expiry/timeout parameters to respect congestion budgets given slack.

\begin{figure}[t]
\centering
\fbox{\begin{minipage}{0.92\linewidth}\centering\small Figure omitted in source-only release. Requested file: \texttt{fig\_slack\_aware\_expiry\_opt\_tv.pdf}.\end{minipage}}
\caption{Archive evidence: SAES trades expiry aggressiveness against congestion distinguishability.}
\label{fig:saes}
\end{figure}

\section{Declared dither and the \texorpdfstring{$\delta$}{delta}-tax for periodic calendars}
A periodic service calendar introduces discretization artifacts: witnesses can exploit periodicity and quantization.
Declared dither (publicly specified jitter) can reduce TV and attenuate spectral witnesses.

\begin{figure}[t]
\centering
\fbox{\begin{minipage}{0.92\linewidth}\centering\small Figure omitted in source-only release. Requested file: \texttt{fig\_dither\_calendar\_tv.pdf}.\end{minipage}}
\caption{Archive evidence: dithered calendars can contract congestion TV compared to rigid calendars.}
\label{fig:dither}
\end{figure}

\subsection{Declared dither as a contractive channel}
A declared-dither mechanism can be viewed as a Markov kernel \(K\) that maps an ``undithered'' calendar-aligned quantity (or witness) \(X\) to an observed quantity \(Y\).
Total variation contracts under such channels.

\begin{lemma}[TV contraction via the Dobrushin coefficient]\label{lem:dobrushin}
Let \(K(\cdot\mid x)\) be a Markov kernel and define its Dobrushin (TV) contraction coefficient
\[
\eta_{\TV}(K) := \sup_{x,x'} \TV\big(K(\cdot\mid x),K(\cdot\mid x')\big)\in[0,1].
\]
Then for any two input laws \(P,Q\) on \(X\),
\[
\TV(PK,QK)\le \eta_{\TV}(K)\,\TV(P,Q)\le \TV(P,Q).
\]
If the same declared-dither kernel is applied \(m\) times (e.g., across ticks or hop-local stages), then
\(\TV(PK^{m},QK^{m})\le \eta_{\TV}(K)^{m}\TV(P,Q)\).
\end{lemma}

\begin{proof}
Standard; see Polyanskiy--Wu for SDPI and contraction coefficients \cite{polyanskiy_wu_sdpi}.
The TV contraction coefficient is also closely related to Dobrushin's ergodicity coefficient for Markov kernels \cite{dobrushin1956,delmoral2003}.
\end{proof}

\begin{remark}[When does dither help?]
The coefficient \(\eta_{\TV}(K)\) is strictly less than \(1\) only when different calendar states induce output distributions with nontrivial overlap at the observer's effective resolution.
This is why PSC-Q couples \((\delta,\alpha)\) to an observer model and why monitoring must confirm that the deployment remained in the assumed timestamp-noise/jitter regime.
\end{remark}

\begin{remark}[Contraction coefficients as ``privacy knobs'']
Recent privacy work uses Dobrushin/TV contraction coefficients as computable summaries of how much a randomized mechanism collapses distinguishability, beyond the local-DP setting \cite{grosse2026}.
In PSC-Q, this motivates exporting \(\eta_{\TV}(K)\) (or a conservative upper bound) as a mechanism-packet-visible knob alongside \((\delta,\alpha)\), so later verifier layers can replay the same declared contraction assumption without restating the mechanism analysis.
\end{remark}


\begin{figure}[t]
\centering
\fbox{\begin{minipage}{0.92\linewidth}\centering\small Figure omitted in source-only release. Requested file: \texttt{fig\_periodic\_calendar\_delta\_tax.pdf}.\end{minipage}}
\caption{New evidence (added in this split archive): quantifying the discretization ``$\delta$-tax'' for periodic calendars.
At fixed utilization $\rho=0.8$ and a 1\% arrival-rate shift, larger $\delta$ increases waiting-time distinguishability, while modest public dither ($\alpha=0.25$) reduces it.}
\label{fig:deltatax}
\end{figure}

\subsection{Robustness to utilization drift}
In deployment, utilization is not constant.
Since the TV tier is slack-sensitive, a PSC-Q knob choice that is safe at one utilization can fail when background load increases.
We therefore add a small robustness sweep varying $\rho$ while holding $(\delta,\alpha)$ fixed.
Fig.~\ref{fig:drift} shows that TV grows rapidly as $\rho$ approaches 1, and that dither can help but does not eliminate the need to monitor and re-tune.

\begin{figure}[t]
\centering
\fbox{\begin{minipage}{0.92\linewidth}\centering\small Figure omitted in source-only release. Requested file: \texttt{fig\_delta\_tax\_utilization\_drift.pdf}.\end{minipage}}
\caption{New evidence (added in this update): utilization-drift robustness for the TV-tier $\delta$-tax experiment.
For several representative $(\delta,\alpha)$ settings, we plot waiting-time TV under a 1\% arrival-rate shift as utilization $\rho$ varies.
This motivates compiling PSC-Q knobs against a utilization range and recording \emph{monitoring evidence} that the deployed regime stayed in-range, while leaving any checker/notarization wrapping of that evidence to later companion layers.}
\label{fig:drift}
\end{figure}



\section{Supplementary knob compilation results (integrated)}
To avoid losing useful engineering-facing content during the transition to a unified archive,
we inline the most deployable ``supplement'' material here in a compact form.

\subsection{Representative knob table}
Table~\ref{tab:pscq_knobs} summarizes one simple PSC-Q knob sweep used by the archive evidence: for several
timestamp-noise levels $\sigma$ and utilization ceilings $\rho_{\max}$, it lists a feasible calendar family and key parameters.
\begin{table}[h]
\centering
\begin{tabular}{r|r|l|r|r|r|r}
$\sigma$ & $\rho_{\max}$ & family & $\delta$ & $\alpha$ & $\beta$ & $1/\delta$ \\
\hline
0.02 & 0.80 & renewal & 2.00 & -- & 0.49 & 0.500\\
0.02 & 0.90 & -- & -- & -- & -- & --\\
0.02 & 0.95 & -- & -- & -- & -- & --\\
0.05 & 0.80 & renewal & 2.00 & -- & 0.49 & 0.500\\
0.05 & 0.90 & -- & -- & -- & -- & --\\
0.05 & 0.95 & -- & -- & -- & -- & --\\
0.10 & 0.80 & renewal & 2.00 & -- & 0.49 & 0.500\\
0.10 & 0.90 & -- & -- & -- & -- & --\\
0.10 & 0.95 & -- & -- & -- & -- & --\\
\end{tabular}
\caption{Representative PSC-Q knob sweep (from the original split-paper evidence).}\label{tab:pscq_knobs}
\end{table}

\subsection{Knob manifest snippets and presets}
PSC-Q can be deployed with a public ``knob manifest'' committing to the chosen calendar family, dither, utilization range, observer assumptions, and monitoring expectations.
The manifest is the stable surface of the PSC-Q mechanism packet; W-Congestion-EQ consumes it as an input rather than redefining its mechanism semantics.

\paragraph{Change control for calendar presets.}
Because knob manifests and calendar families are public artifacts, updates should run under the same disagreement-gated change-control discipline used for certified routing artifacts: log a proposed preset update, attach monitoring residuals (e.g., drift in slot-index recovery rates or utilization), and accept only if the disagreement gate is satisfied and the updated manifest is transparency-anchored \cite{series:recert,anondht_release_receipts}.

\paragraph{Minimal manifest skeleton.}
\begin{lstlisting}[basicstyle=\ttfamily\small]
\begin{quote}\small
\texttt{version}: \texttt{pscq-knobs-v0.1}\\
\texttt{calendar.family}: \texttt{renewal\_uniform}\\
\texttt{calendar.delta\_ms}: \texttt{2}\\
\texttt{calendar.beta\_frac}: \texttt{0.49}\\
\texttt{utilization}: [0.7, 0.8]\\
\texttt{budgets}: (eps\_TV=0.05, eps\_spectral=0.15)\\
\texttt{knob\_table\_sha256}: \texttt{90a6a2c886069871...}\\
\end{quote}
\end{lstlisting}

\paragraph{Preset examples (low / default / high noise).}
\begin{lstlisting}[basicstyle=\ttfamily\small]
\begin{quote}\small
\texttt{version}: \texttt{pscq-knobs-v0.1}\\
\texttt{calendar.family}: \texttt{renewal\_uniform}\\
\texttt{calendar.delta\_ms}: \texttt{2}\\
\texttt{calendar.beta\_frac}: \texttt{0.49}\\
\texttt{utilization}: [0.7, 0.8]\\
\texttt{budgets}: (eps\_TV=0.05, eps\_spectral=0.15)\\
\texttt{knob\_table\_sha256}: \texttt{90a6a2c886069871...}\\
\end{quote}
\end{lstlisting}


\smallskip
\noindent Other presets (e.g., higher observation noise or tighter utilization ceilings) use the same schema with different $(\delta,\alpha,\beta)$ and different \texttt{knob\_table\_sha256} values. They are omitted here to keep the paper tight; the full preset table and reproduction scripts are available in the on-request artifact bundle.

\subsection{Two short notes: utilization-robust compilation and slot-index recovery}
\paragraph{Interval-robust compilation against utilization drift.}
A practical compiler target is a utilization range $\rho\in[\rho_{\min},\rho_{\max}]$ rather than a point value.
A minimal robustification is: pick $(\delta,\alpha)$ to satisfy the Congestion-EQ TV contract at $\rho_{\max}$,
and publish $\rho_{\max}$ as part of the manifest so that monitoring can trigger recertification when the regime changes \cite{series:recert}.
This is the ``range compilation'' view used by the utilization-drift sweep in Fig.~\ref{fig:drift}.

\paragraph{De-jittering and slot-index recovery (observer model).}
Let the calendar tick width be $\delta$ and let declared dither add uniform jitter in $[-\alpha\delta,\alpha\delta]$.
If the observer has effective timestamp noise $\sigma_{\mathrm{obs}}$ (including instrumentation error and network jitter),
then a crude sufficient condition for frequent slot-index ambiguity is
\[
(1/2-\alpha)\,\delta \;\lesssim\; \sigma_{\mathrm{obs}},
\]
since adjacent ticks overlap at the observer's resolution.
Conversely, if $\sigma_{\mathrm{obs}}\ll (1/2-\alpha)\delta$, an observer can often ``snap'' timestamps back to the nearest tick
and recover slot indices, reintroducing periodic witnesses. This motivates treating $(\delta,\alpha)$ as certificate-visible knobs
and monitoring for drift in $\sigma_{\mathrm{obs}}$.


\section{Microstudies summary (plots on request)}
\label{sec:microstudies-summary}
The core PSC-Q argument is mechanism-level; nevertheless, a few deployment-facing microstudies are worth stating explicitly because they recur across systems.
Figures and raw tables are omitted in this source-only archive but are available on request.
\begin{itemize}[leftmargin=*]
\item \textbf{Additive padding can worsen congestion witnesses:} when additive cover pushes utilization into the high-delay knee, small secret-correlated load differences can create large shifts in waiting-time/RTT laws (Congestion-EQ tier). PSC-Q therefore prefers slack-harvesting substitution.
\item \textbf{Partial idle-filling is often enough:} filling only a fraction of idle calendar slots can still materially reduce trace-tier distinguishability while preserving slack; the right fraction depends on the congestion budget and the downstream resource model.
\item \textbf{Timestamp resolution matters:} if observers can recover slot indices from high-resolution timestamps, periodic/spectral structure returns. Declared dither should be sized so that slot ambiguity dominates effective observer timestamp noise.
\end{itemize}

\section{Practical guidance for Anonymous-DHT compilers}
\begin{itemize}[leftmargin=1.2em]
\item Prefer substitution/idle-filling to additive padding when the goal is congestion privacy.
\item Provision explicit slack targets; do not treat ``cover rate'' as an independent knob.
\item Treat $\delta$ and dither as first-class compiler outputs and include them in the declared PSC-Q mechanism packet / knob manifest that later verifier layers consume.
\item If a periodic/spectral tier is in scope, choose \(\alpha\) and \(\delta\) so that slot ambiguity exceeds timestamp resolution.
  A crude rule-of-thumb from the slot-index recovery note above is to target \((1/2-\alpha)\,\delta \lesssim \sigma_{\mathrm{obs}}\), where \(\sigma_{\mathrm{obs}}\) is the observer's effective timestamp noise/jitter.
\item Evaluate a witness family across hops (RTT time series, queue delay, loss/ECN) to avoid bottleneck shift.
\end{itemize}

\section{Takeaways}
\begin{itemize}[leftmargin=*]
\item \textbf{Substitution beats additive for congestion privacy:} use idle-filling to reduce trace witnesses without shrinking slack.
\item \textbf{Calendars are public interfaces:} when service opportunities are public and auditable, the remaining secrecy surface is pushed into controlled degrees of freedom that can be budgeted.
\item \textbf{Dither is a first-class knob:} declared dither can contract TV-tier witnesses, but periodic/spectral tiers require explicit attention to timestamp resolution and slot recovery.
\item \textbf{Treat utilization as part of the privacy contract:} congestion witnesses are highly nonlinear near the knee; publish slack targets and monitor drift.
\end{itemize}

\section{Limitations and open questions}
Declared dither contracts the \emph{TV tier} for waiting-time witnesses, but it does not automatically protect against all witnesses.
Periodic structure is routinely detectable in encrypted-traffic time series using autocorrelation and Lomb--Scargle periodograms \cite{lomb1976,scargle1982,koumar2022periodic,assadhan2014,periodicity2025}, so a spectral tier is a realistic concern for overlays whose calendars induce periodicity. For a formal spectral-tier treatment in the series (including optimal distinguishing bounds under delegation), see \cite{series:spectral}.
As the slot-index recovery note above shows, an observer with high-resolution timestamps can often de-jitter and recover slot indices, regaining periodic structure.
A full treatment therefore benefits from a tiered view: (i) TV-tier accounting (contract paper), plus (ii) explicit periodic/spectral-tier goals, and (iii) monitoring for drift and unexpected witnesses (assurance paper).


\section{Artifacts and reproducibility}
This source-only release is intentionally small: it publishes the referee-facing arguments as a single \\LaTeX\\ file.
Figures and full simulation artifacts used to generate the evidence are available on request as an external archive.
To keep claims auditable, PSC-Q stops at the mechanism packet itself; W-Congestion-EQ specifies the congestion-specific verifier-bundle fields that later bind that packet to evidence digests, replay profiles, and checker outputs.


\begin{thebibliography}{99}
\bibitem{series:congeq}
Anonymous.
\newblock \emph{Congestion-EQ: Total-Variation Contracts for Queue/RTT Leakage in Anonymous Overlays and Anonymous DHT Lookups}.
\newblock Contract paper in this unified archive.

\bibitem{series:schedcompiler}
Anonymous.
\newblock \emph{Scheduling and Compiler Techniques for Metadata-Hiding Overlay Lookups}.
\newblock Draft v0.67 (2026-01-21). Series paper.
\newblock Published. Wiki: \texttt{[[2026.01.22 - Mathematics: Scheduling and Compiler Techniques for Metadata-Hiding Overlay Lookups]]}.

\bibitem{series:knapsacktransport}
Anonymous.
\newblock \emph{Optimal Poset-Feasible Padding: Knapsack, Transport, and Dual Certificates for TV Privacy}.
\newblock January 2026. Published. Wiki: \texttt{[[2026.01.29 - Mathematics: Optimal Poset-Feasible Padding: Knapsack, Transport, and Dual Certificates for TV Privacy]]}.

\bibitem{series:spectral}
Anonymous.
\newblock \emph{Spectral Anonymity and Optimal Distinguishing Bounds for Random-Walk Delegation in (Possibly Directed) Overlays}.
\newblock January 2026. Series paper.
\newblock Published. Wiki: \texttt{[[2026.01.23 - Mathematics: Spectral Anonymity and Optimal Distinguishing Bounds for Random-Walk Delegation in (Possibly Directed) Overlays]]}.

\bibitem{series:wcongeq}
Anonymous.
\newblock \emph{W-Congestion-EQ: Verifiable Congestion Budgets --- Congestion-Specific Verifier Bundles for Anonymous Overlays}.
\newblock Assurance companion in this split archive.

\bibitem{series:pscqaddendum}
Anonymous.
\newblock \emph{PSC-Q (Addendum): Supplementary Evidence and Operator Microstudies}.
\newblock Congestion series Paper~2A, 2026. (This archive.)

\bibitem{series:recert}
Anonymous.
\newblock \emph{Disagreement-Gated Recertification for Certified Anonymous-DHT Routing Artifacts}.
\newblock Certified-series companion in this unified archive (change control and drift gates).

\bibitem{anondht_release_receipts}
Anonymous.
\newblock \emph{Privacy as a release property: ledger bundles and receipts for anonymous DHT deployments}.
\newblock Release-and-destination companion in this unified archive (receipt schema, transparency anchoring).

\bibitem{series:evalnotary}
Anonymous.
\newblock \emph{Notarized Anonymity Certificates Under Retries: Grinding-Resistant Evaluation via Transparency Logs, Public Randomness, and $\delta$-Spending}.
\newblock Evaluation series Paper~1 in this unified archive.

\bibitem{series:opb}
Anonymous.
\newblock \emph{Auditable Trace Privacy for Anonymous DHTs: Runtime Plan Pinning, PTL Checks, and Epoch Conformance Receipts}.
\newblock Operational series Paper~B in this unified archive.

\bibitem{torproposal254}
Tor Project.
\newblock Proposal 254: Padding Negotiation.
\newblock Tor design proposal, 2015 (closed; implemented as \texttt{padding-spec}).
\newblock \url{https://spec.torproject.org/proposals/254-padding-negotiation.html}.

\bibitem{torpaddingspec}
M.~Perry and G.~Kadianakis.
\newblock Tor Padding Specification.
\newblock Tor protocol specification, 2019--.
\newblock \url{https://spec.torproject.org/padding-spec/}.

\bibitem{dyer2012peekaboo}
K.~P. Dyer, S.~E. Coull, T.~Ristenpart, and T.~Shrimpton.
\newblock Peek-a-Boo, I Still See You: Why Efficient Traffic Analysis Countermeasures Fail.
\newblock In \emph{IEEE Symposium on Security and Privacy}, 2012.
\newblock DOI: 10.1109/SP.2012.28.

\bibitem{csbuflo}
X.~Cai, R.~Nithyanand, R.~Johnson, and I.~Goldberg.
\newblock CS-BuFLO: A Congestion Sensitive Website Fingerprinting Defense.
\newblock In \emph{WPES}, 2014.
\newblock DOI: 10.1145/2665943.2665949.

\bibitem{tor324}
Tor Project.
\newblock Proposal 324: RTT Congestion Control.
\newblock Tor design proposal, 2020--.
\newblock \url{https://spec.torproject.org/proposals/324-rtt-congestion-control.html}.

\bibitem{torcc2022}
Tor Project.
\newblock Congestion Control Arrives in Tor 0.4.7-stable!
\newblock Tor Project blog, May 4, 2022.
\newblock \url{https://blog.torproject.org/congestion-contrl-047/}.

\bibitem{defenestrator}
M.~AlSabah, K.~Bauer, I.~Goldberg, D.~Grunwald, D.~McCoy, S.~Savage, and G.~M. Voelker.
\newblock DefenestraTor: Throwing Out Windows in Tor.
\newblock In \emph{PETS}, 2011.
\newblock DOI: 10.1007/978-3-642-22263-4\_8.

\bibitem{juarez_ap}
M.~Ju\'arez, M.~Imani, M.~Perry, C.~D\'iaz, and M.~Wright.
\newblock Toward an efficient website fingerprinting defense.
\newblock In \emph{ESORICS}, 2016.\newblock DOI: 10.1007/978-3-319-45744-4\_2.\newblock \url{https://nymity.ch/tor-dns/bibliography/pdf/Juarez2016a.pdf}.

\bibitem{witwer2022}
E.~Witwer, J.~Holland, and N.~Hopper.
\newblock Padding-only defenses add delay in Tor.
\newblock In \emph{Workshop on Privacy in the Electronic Society (WPES)}, 2022.
\newblock DOI: 10.1145/3559613.3563207. Also arXiv:2208.02917.
\newblock \url{https://dl.acm.org/doi/10.1145/3559613.3563207}.

\bibitem{pulls_interspace}
T.~Pulls.
\newblock Towards effective and efficient padding machines for Tor.
\newblock arXiv:2011.13471, 2020.
\newblock \url{https://arxiv.org/abs/2011.13471}.

\bibitem{shaperd2025}
S.~Wilson, S.~Tian, and S.~Kamali.
\newblock Extended Abstract: Shaperd: Easily Adoptable Real-Time Traffic Shaper for Fully Encrypted Protocols.
\newblock Free and Open Communications on the Internet (FOCI), 2025.
\newblock \url{https://petsymposium.org/foci/2025/foci-2025-0010.pdf}.

\bibitem{lomb1976}
N.~R. Lomb.
\newblock Least-squares frequency analysis of unequally spaced data.
\newblock \emph{Astrophysics and Space Science}, 39:447--462, 1976.

\bibitem{scargle1982}
J.~D. Scargle.
\newblock Studies in astronomical time series analysis. II. Statistical aspects of spectral analysis of unevenly spaced data.
\newblock \emph{The Astrophysical Journal}, 263:835--853, 1982.

\bibitem{koumar2022periodic}
J.~Koumar and T.~\v{C}ejka.
\newblock Network traffic classification based on periodic behavior detection.
\newblock In \emph{18th International Conference on Network and Service Management (CNSM)}, 2022.
\newblock \url{https://dl.ifip.org/db/conf/cnsm/cnsm2022/52.pdf}.

\bibitem{assadhan2014}
B.~AsSadhan and J.~M.~F. Moura.
\newblock An efficient method to detect periodic behavior in botnet traffic by analyzing control plane traffic.
\newblock \emph{Journal of Advanced Research}, 5(4):435--448, 2014.
\newblock DOI: 10.1016/j.jare.2013.11.005.

\bibitem{periodicity2025}
D.~O\v{s}kera, J.~Koumar, A.~Pokorn\'a, K.~Je\v{r}\'abek, and T.~\v{C}ejka.
\newblock Botnet detection through periodic patterns in command-and-control network traffic.
\newblock In \emph{21st International Conference on Network and Service Management (CNSM)}, 2025.
\newblock DOI: 10.23919/CNSM67658.2025.11297465.
\newblock \url{https://dl.ifip.org/db/conf/cnsm/cnsm2025/1571190192.pdf}.


\bibitem{agat2000}
J.~Agat.
\newblock Transforming out timing leaks.
\newblock In \emph{Proceedings of the 27th ACM SIGPLAN-SIGACT Symposium on Principles of Programming Languages (POPL)}, 2000, pp.~40--53.
\newblock DOI: 10.1145/325694.325702.

\bibitem{mishra_fallthrough2025}
A.~Mishra and A.~Bichhawat.
\newblock Fall-Through Semantics for Mitigating Timing-Based Side Channel Leaks.
\newblock In \emph{FSTTCS} (LIPIcs), 2025.
\newblock DOI: 10.4230/LIPIcs.FSTTCS.2025.44.
\newblock \url{https://drops.dagstuhl.de/entities/document/10.4230/LIPIcs.FSTTCS.2025.44}.


\bibitem{falkner_imc2007}
J.~Falkner, M.~Piatek, J.~P.~John, A.~Krishnamurthy, and T.~Anderson.
\newblock Profiling a Million User DHT.
\newblock In \emph{ACM SIGCOMM Internet Measurement Conference (IMC)}, 2007.
\newblock DOI: 10.1145/1298306.1298325.
\newblock \url{https://homes.cs.washington.edu/~arvind/papers/dht.pdf}.

\bibitem{ipfs_kad_dht_spec}
IPFS Standards.
\newblock IPFS Kademlia DHT Specification (Concurrency parameter \(\alpha\) and client lookup behavior).
\newblock 20 November 2025.
\newblock \url{https://specs.ipfs.tech/routing/kad-dht/}.

\bibitem{stutzbach_rejaie2006}
D.~Stutzbach and R.~Rejaie.
\newblock Improving Lookup Performance over a Widely-Deployed DHT.
\newblock In \emph{IEEE INFOCOM}, 2006.
\newblock DOI: 10.1109/INFOCOM.2006.329.
\newblock \url{https://www.cs.uoregon.edu/Reports/TR-2005-005.pdf}.

\bibitem{cortes_kademlia_das2024}
M.~Cortes-Goicoechea, C.~Kiraly, D.~Ryajov, J.~L. Munoz-Tapia, and L.~Bautista-Gomez.
\newblock Scalability limitations of Kademlia DHTs when enabling Data Availability Sampling in Ethereum.
\newblock arXiv:2402.09993, 2024.
\newblock \url{https://arxiv.org/abs/2402.09993}.

\bibitem{polyanskiy_wu_sdpi}
Y.~Polyanskiy and Y.~Wu.
\newblock Strong data-processing inequalities for channels and Bayesian networks.
\newblock arXiv:1508.06025, 2015.
\newblock \url{https://arxiv.org/abs/1508.06025}.

\bibitem{dobrushin1956}
R.~L. Dobrushin.
\newblock Central limit theorem for nonstationary Markov chains.
\newblock \emph{Theory of Probability and Its Applications}, 1(1):72--89, 1956.
\newblock \url{https://www.mathnet.ru/eng/tvp4987}.

\bibitem{delmoral2003}
P.~Del~Moral, M.~Ledoux, and L.~Miclo.
\newblock On contraction properties of Markov kernels.
\newblock \emph{Probability Theory and Related Fields}, 126:395--420, 2003.
\newblock DOI: 10.1007/s00440-003-0270-6. \url{https://perso.math.univ-toulouse.fr/miclo/files/2012/04/contraction.pdf}.

\bibitem{grosse2026}
L.~Grosse, S.~Saeidian, T.~J. Oechtering, and M.~Skoglund.
\newblock Dobrushin Coefficients of Private Mechanisms Beyond Local Differential Privacy.
\newblock arXiv:2601.09498, 2026.
\newblock \url{https://arxiv.org/abs/2601.09498}.


\bibitem{series:artifactpolicy}
Anonymous.
\newblock Artifact policy (source-only archive; artifacts available on request).
\newblock Series note, 2026.
\end{thebibliography}

\end{document}